% Generated by Listas/make_statements.py from pratica-lecture-02.md, item 19. Do not edit.
\textbf{Source.} Mas-Colell, Whinston \& Green (1995), \emph{Microeconomic Theory}, Ch. 8 Exercises, Exercise 8.D.3\^{}B, p.~263 {[}PDF p.~142{]}. Answer in the source: none in the book (the MWG solutions manual covers Chapter 8). Reproduced verbatim.



\begin{quote}

Consider a first-price sealed-bid auction of an object with two bidders. Each bidder \(i\)'s valuation of the object is \(v_i\), which is known to both bidders. The auction rules are that each player submits a bid in a sealed envelope. The envelopes are then opened, and the bidder who has submitted the highest bid gets the object and pays the auctioneer the amount of his bid. If the bidders submit the same bid, each gets the object with probability \(\frac{1}{2}\). Bids must be in dollar multiples (assume that valuations are also).



\textbf{(a)} Are any strategies strictly dominated?



\textbf{(b)} Are any strategies weakly dominated?



\textbf{(c)} Is there a Nash equilibrium? What is it? Is it unique?

\end{quote}

